% !TEX encoding = UTF-8 Unicode
\section{Fachliches Umfeld}
\subsection{Der Sampleblock}
Ein digitales Audiosignal besteht aus einer Menge einzelner Werte, den Samples. Der Sampleblock stellt die Menge an Samples dar, die von einem VST-Plugin auf einmal bearbeitet werden kann. Die Größe des Sampleblockes ist Ausschlaggebend für die Performance und die Latenz\footnote{die Verzögerung zwischen Signalein- und ausgabe.} einer DAW Umgebung. Je niedriger die Blockgröße desto niedriger die Latenz, desto höher die CPU Auslastung.
\subsection{\label{Prozess}Der Prozess}
In der Terminologie des VST-SDK, ist mit Prozess der Teil eines Plugins gemeint, der für die Bearbeitung aller Pluginein- und ausgänge verantwortlich ist.
\\\\
Im VST2.x SDK ist dies eine Methode namens ''processReplacing''. Ihr werden als Parameter die Eingangs- und ausgangssampleblöcke übergeben. Die Anzahl der Sampleblöcke ist abhängig von der Anzahl der Pluginein- sowie ausgangskanäle. So werden zum Beispiel einem Stereoplugin zwei Eingangsampleblöcke übergeben, jeweils ein Block für ''Rechts'' und ein Block für ''Links''.
\\\\
Der Ablauf eines Prozesses kann wie folgt Skizziert werden:
\begin{itemize}
	\item die DAW füllt die Plugin-Eingangssampleblöcke mit den Daten eines Audiokanals
	\item die DAW ruft die ''processReplacing'' Methode des Plugins auf.
	\item das Plugin verarbeitet die Eingangssampleblöcke und schreibt die Ergebnisse in die entsprechenden Ausgangssampleblöcke
	\item die DAW verarbeitet die Ausgangssampleblöcke weiter
\end{itemize}
Dieser Vorgang wird, während eine DAW im ''Abspielmodus'', ist ständig wiederholt.
